\section{The Grounding Gap}
\label{sec:grounding-gap}

There are two distinct kinds of reliability a language model can have.
\emph{Empirical reliability} means the model produces correct outputs most
of the time on a given distribution, as measured by benchmarks or human
evaluation. \emph{Formal reliability} means the model's outputs are
verifiable as correct relative to a specified knowledge base and ontology.

Current LLMs have empirical reliability. Grounded BP systems --- transformers
operating over a declared factor graph with a verifier --- have formal
reliability in restricted settings: outputs can be checked against the
declared knowledge base and failures traced to specific premises. The gap
between these two is the grounding gap.

\subsection*{Why Scaling Does Not Close It}

Scaling improves empirical performance on many benchmark distributions, but
the relationship between scale and reliability is task-dependent ---
calibration can worsen and failure modes can shift. More fundamentally,
scaling does not create a declared ontology, a finite concept space, or a
verifier. A 405B parameter model with no grounding structure cannot provide
formal correctness guarantees, regardless of its benchmark performance.

The grounding gap is not a quantitative gap that more parameters will close.
It is a structural gap between two different kinds of system.

\subsection*{What Would Close It}

Closing the grounding gap requires building systems that combine neural
inference (the transformer's BP computation) with symbolic grounding (a
declared ontology and verifier). This is the QBBN program. Whether it can
be scaled to general-purpose language understanding is the central open
problem.

\subsection*{Epistemic Status}

\begin{center}
\begin{tabular}{ll}
\toprule
Claim & Status \\
\midrule
Two kinds of reliability are distinct & Conceptual distinction \\
Grounded BP systems have formal verifiability & Direct consequence \\
Scaling reliably improves all reliability metrics & Explicitly not claimed \\
Scaling does not close the grounding gap & Mechanistic interpretation \\
QBBN hybrid closes the gap & Architectural proposal \\
\bottomrule
\end{tabular}
\end{center}